\chapter{Absztrakt}
\label{ch:abstract}

A dolgozat célja egy nagyméretű, saját készítésű képadatbázis létrehozása és ezen mély neurális hálózatok osztályozási teljesítményének vizsgálata volt. 
Az SBphotothemes180k nevű adatbázis 180 750 egyedi, címkézett fényképet tartalmaz, kétféle annotációval: az ImageNet1K-ből szűrt 114, valamint a Places365-ből származó 121 kategóriával. Az osztályeloszlás erősen torzított, természetes adatok esetén tipikus, long-tail eloszlást mutatott (aszimmetria(skewness): 8,4; szórás: ~3200), de a nagy adatmennyiség és a szortírozási metodika révén kellően nagyra adódott az átlagos kategóriaméret is (1500). 
Az adatokon háromfázisú tesztelést végeztem hét népszerű, de különböző felépítésű neurális hálózaton (pl. ResNet50, EfficientNet-B0, ViT-B/16), különböző tanítóhalmaz-méretekkel. 
Az első lépés a tanítás nélküli modellek tesztelése volt a beépített súlyokkal, ezt követte a teljes tanítás, majd végül az iteratívan növekvő tanítóhalmazokon történő tesztelés. A modellek közül összességében az EfficientNet-B0 és a ConvNeXt-Tiny emelkedtek ki. A teljes adathalmazon tanított modellek közül az EfficientNet-B0 és a NASNet mutatták a legjobb teljesítményt (Top-1 pontosság ~76\%, Top-5 ~94\%). Az utolsó fázisban a fokozatosan növelt tanítóhalmazokon mért teljesítményt vizsgáltam, hogy látható legyen a tanítóhalmaz méretének hatása.
Ezen túl saját, erőforrástakarékos osztályozó modellt is fejlesztettem ResNet18-alapon, különféle fejrészekkel. Az egyszerű egyrétegű modell hasonló pontosságot ért el, mint az új modern fejrész, de az új modell számítási igénye jóval kisebb és pontossága összemérhető a nagy modellekével, így mobil eszközökön való alkalmazásra is alkalmas lehet. 
A nagyméretű adathalmaz, a többféle osztályozás, a nagy mennyiségű mért adat és a saját modellek további vizsgálatok alapját képezhetik a jövőben.